% GENERATED FILE. Edit canonical lessons, then run npm run generate:books.
% canonical-source-hash: fnv1a64:8e36c99dfe81267e
% canonical-lessons: PA-C232-sajauna, PA-C232-sametna, PA-C232-yad, PA-C232-uddharna, PA-C232-tunnna

\chapter{Putting Things Away}
\label{ch:pa-a2-putting-things-away}
\hlchaptermodality{232}

\begin{chapteropening}
\textbf{By the end of this chapter:} \emph{I can say how I keep my things in order.}

Complete the last lesson of chapter 232: I can say how I keep my things in order.
\end{chapteropening}

\section[\texorpdfstring{\pa{ਸਜਾਉਣਾ} (sajāuṇā)}{ਸਜਾਉਣਾ (sajāuṇā)}]{\pa{ਸਜਾਉਣਾ} (sajāuṇā) — to arrange neatly; to decorate}

\label{lesson:PA-C232-sajauna}



\begin{quote}
Before the new one: say the Punjabi for a hair tassel, then the Punjabi for a lock of hair.
\end{quote}

\subsection*{You'll want to know: \pa{ਸਜਾਉਣਾ}}
\textbf{\pa{ਸਜਾਉਣਾ}} — \emph{sajāuṇā} — \textquotedblleft{}to arrange neatly; to decorate\textquotedblright{}.

\begin{quote}
\pa{ਮੈਂ} \pa{ਅਲਮਾਰੀ} \pa{ਵਿੱਚ} \pa{ਕਿਤਾਬਾਂ} \pa{ਸਜਾਉਂਦੀ} \pa{ਹਾਂ।}
\end{quote}

\begin{quote}
\emph{maiṁ almārī vich kitābāṁ sajāundī hāṁ.}
\end{quote}

\begin{quote}
I arrange the books neatly in the cupboard.
\end{quote}

Past (I …): \pa{ਮੈਂ} \pa{ਕਿਤਾਬਾਂ} \pa{ਸਜਾਈਆਂ}, \emph{maiṁ kitābāṁ sajāīāṁ}.

\begin{cousinweb}
\textbf{\pa{ਸਜਾਉਣਾ}} (\emph{sajāuṇā}) is inherited from Sanskrit \emph{sajjayati}, \textquotedblleft{}equips, prepares\textquotedblright{}; CDIAL lists Punjabi \emph{sajāuṇā}, \textquotedblleft{}to prepare\textquotedblright{} and Lahnda \emph{sajāvaṇ}, \textquotedblleft{}to adorn\textquotedblright{} (CDIAL 13091), but does not give the sense \textquotedblleft{}arrange neatly\textquotedblright{}.
\end{cousinweb}

\begin{grammarlens}[title={What you've built}]
One more word for saying what you do.
\end{grammarlens}

\subsection*{Your turn}
\begin{itemize}
\raggedright
  \item \emph{Say it:} \emph{sajāuṇā}
  \item \emph{Say it:} \emph{sajāuṇā}, once more
  \item \emph{Say it:} \emph{maiṁ almārī vich kitābāṁ sajāundī hāṁ.} — I arrange the books neatly in the cupboard.
  \item \emph{Recall:} say the Punjabi for a hair tassel, then the Punjabi for a lock of hair, then say \emph{sajāuṇā} again
\end{itemize}

\begin{tcolorbox}[breakable,colback=teal!4,colframe=teal!35!black,title={Before you move on}]
What does \pa{ਸਜਾਉਣਾ} mean? (\textquotedblleft{}To arrange neatly; to decorate\textquotedblright{}.) Say it once more.
\end{tcolorbox}
\section[\texorpdfstring{\pa{ਸਮੇਟਣਾ} (sameṭṇā)}{ਸਮੇਟਣਾ (sameṭṇā)}]{\pa{ਸਮੇਟਣਾ} (sameṭṇā) — to gather up; to roll up}

\label{lesson:PA-C232-sametna}



\begin{quote}
Before the new one: say the Punjabi for a lock of hair, then the Punjabi for to arrange neatly.
\end{quote}

\subsection*{You'll want to know: \pa{ਸਮੇਟਣਾ}}
\textbf{\pa{ਸਮੇਟਣਾ}} — \emph{sameṭṇā} — \textquotedblleft{}to gather up; to roll up\textquotedblright{}.

\begin{quote}
\pa{ਸਵੇਰੇ} \pa{ਮੈਂ} \pa{ਰਜਾਈ} \pa{ਸਮੇਟਦਾ} \pa{ਹਾਂ।}
\end{quote}

\begin{quote}
\emph{savere maiṁ rajāī sameṭdā hāṁ.}
\end{quote}

\begin{quote}
In the morning I roll up the quilt.
\end{quote}

Past (I …): \pa{ਮੈਂ} \pa{ਕੱਪੜੇ} \pa{ਸਮੇਟੇ}, \emph{maiṁ kappṛe sameṭe}.

\begin{cousinweb}
\textbf{\pa{ਸਮੇਟਣਾ}} (\emph{sameṭṇā}) is inherited from Sanskrit \emph{saṁvēṣṭayati}, \textquotedblleft{}wraps up\textquotedblright{}; CDIAL lists Punjabi \emph{sameṭṇā}, \textquotedblleft{}to roll up\textquotedblright{} (CDIAL 13026).
\end{cousinweb}

\begin{grammarlens}[title={What you've built}]
One more word for saying what you do.
\end{grammarlens}

\subsection*{Your turn}
\begin{itemize}
\raggedright
  \item \emph{Say it:} \emph{sameṭṇā}
  \item \emph{Say it:} \emph{sameṭṇā}, once more
  \item \emph{Say it:} \emph{savere maiṁ rajāī sameṭdā hāṁ.} — In the morning I roll up the quilt.
  \item \emph{Recall:} say the Punjabi for a lock of hair, then the Punjabi for to arrange neatly, then say \emph{sameṭṇā} again
  \item \emph{Write it:} copy \textbf{\pa{ਸਮੇਟਣਾ}} once, checking the whole word
\end{itemize}

\begin{tcolorbox}[breakable,colback=teal!4,colframe=teal!35!black,title={Before you move on}]
What does \pa{ਸਮੇਟਣਾ} mean? (\textquotedblleft{}To gather up; to roll up\textquotedblright{}.) Say it once more.
\end{tcolorbox}
\section[\texorpdfstring{\pa{ਯਾਦ} (yād)}{ਯਾਦ (yād)}]{\pa{ਯਾਦ} (yād) — memory; remembering}

\label{lesson:PA-C232-yad}



\begin{quote}
Before the new one: say the Punjabi for to arrange neatly, then the Punjabi for to gather up.
\end{quote}

\subsection*{You'll want to know: \pa{ਯਾਦ}}
\textbf{\pa{ਯਾਦ}} — \emph{yād} — \textquotedblleft{}memory; remembering\textquotedblright{}.

\begin{quote}
\pa{ਮੈਨੂੰ} \pa{ਤੁਹਾਡਾ} \pa{ਨਾਂ} \pa{ਯਾਦ} \pa{ਹੈ।}
\end{quote}

\begin{quote}
\emph{mainū̃ tuhāḍā nāṁ yād hai.}
\end{quote}

\begin{quote}
I remember your name.
\end{quote}

A feminine noun.

\begin{cousinweb}
\textbf{\pa{ਯਾਦ}} (\emph{yād}) is borrowed: it is usually traced to Persian \emph{yād}, \textquotedblleft{}memory, remembrance\textquotedblright{}.
\end{cousinweb}

\begin{grammarlens}[title={What you've built}]
One more word for saying what you do.
\end{grammarlens}

\subsection*{Your turn}
\begin{itemize}
\raggedright
  \item \emph{Say it:} \emph{yād}
  \item \emph{Say it:} \emph{yād}, once more
  \item \emph{Say it:} \emph{mainū̃ tuhāḍā nāṁ yād hai.} — I remember your name.
  \item \emph{Recall:} say the Punjabi for to arrange neatly, then the Punjabi for to gather up, then say \emph{yād} again
  \item \emph{Write it:} copy \textbf{\pa{ਯਾਦ}} once, checking the whole word
\end{itemize}

\begin{tcolorbox}[breakable,colback=teal!4,colframe=teal!35!black,title={Before you move on}]
What does \pa{ਯਾਦ} mean? (\textquotedblleft{}Memory; remembering\textquotedblright{}.) Say it once more.
\end{tcolorbox}
\section[\texorpdfstring{\pa{ਉੱਧੜਨਾ} (uddhaṛnā)}{ਉੱਧੜਨਾ (uddhaṛnā)}]{\pa{ਉੱਧੜਨਾ} (uddhaṛnā) — to come unstitched}

\label{lesson:PA-C232-uddharna}



\begin{quote}
Before the new one: say the Punjabi for to gather up, then the Punjabi for memory.
\end{quote}

\subsection*{You'll want to know: \pa{ਉੱਧੜਨਾ}}
\textbf{\pa{ਉੱਧੜਨਾ}} — \emph{uddhaṛnā} — \textquotedblleft{}to come unstitched\textquotedblright{}.

\begin{quote}
\pa{ਜੇ} \pa{ਧਾਗਾ} \pa{ਕਮਜ਼ੋਰ} \pa{ਹੈ} \pa{ਤਾਂ} \pa{ਜੇਬ} \pa{ਉੱਧੜੇਗੀ।}
\end{quote}

\begin{quote}
\emph{je dhāgā kamzor hai tāṁ jeb uddhṛegī.}
\end{quote}

\begin{quote}
If the thread is weak, the pocket will come unstitched.
\end{quote}

Past (it …): \pa{ਉਹ} \pa{ਜੇਬ} \pa{ਉੱਧੜੀ}, \emph{oh jeb uddhṛī}.

\begin{cousinweb}
\textbf{\pa{ਉੱਧੜਨਾ}} (\emph{uddhaṛnā}) is inherited from Sanskrit \emph{uddharati}, \textquotedblleft{}takes out\textquotedblright{}; CDIAL lists Punjabi \emph{uddhaṛṇā}, \textquotedblleft{}to be ripped\textquotedblright{} (CDIAL 2009).
\end{cousinweb}

\begin{grammarlens}[title={What you've built}]
One more word for saying what you do.
\end{grammarlens}

\subsection*{Your turn}
\begin{itemize}
\raggedright
  \item \emph{Say it:} \emph{uddhaṛnā}
  \item \emph{Say it:} \emph{uddhaṛnā}, once more
  \item \emph{Say it:} \emph{je dhāgā kamzor hai tāṁ jeb uddhṛegī.} — If the thread is weak, the pocket will come unstitched.
  \item \emph{Recall:} say the Punjabi for to gather up, then the Punjabi for memory, then say \emph{uddhaṛnā} again
  \item \emph{Write it:} copy \textbf{\pa{ਉੱਧੜਨਾ}} once, checking the whole word
\end{itemize}

\begin{tcolorbox}[breakable,colback=teal!4,colframe=teal!35!black,title={Before you move on}]
What does \pa{ਉੱਧੜਨਾ} mean? (\textquotedblleft{}To come unstitched\textquotedblright{}.) Say it once more.
\end{tcolorbox}
\section[\texorpdfstring{\pa{ਤੁੰਨਣਾ} (tunnṇā)}{ਤੁੰਨਣਾ (tunnṇā)}]{\pa{ਤੁੰਨਣਾ} (tunnṇā) — to stuff; to cram}

\label{lesson:PA-C232-tunnna}



\begin{quote}
Before the new one: say the Punjabi for memory, then the Punjabi for to come unstitched.
\end{quote}

\subsection*{You'll want to know: \pa{ਤੁੰਨਣਾ}}
\textbf{\pa{ਤੁੰਨਣਾ}} — \emph{tunnṇā} — \textquotedblleft{}to stuff; to cram\textquotedblright{}.

\begin{quote}
\pa{ਮੈਂ} \pa{ਕੱਪੜੇ} \pa{ਬਸਤੇ} \pa{ਵਿੱਚ} \pa{ਤੁੰਨਦਾ} \pa{ਹਾਂ।}
\end{quote}

\begin{quote}
\emph{maiṁ kappṛe baste vich tunndā hāṁ.}
\end{quote}

\begin{quote}
I stuff my clothes into the bag.
\end{quote}

Past (I …): \pa{ਮੈਂ} \pa{ਕੱਪੜੇ} \pa{ਤੁੰਨੇ}, \emph{maiṁ kappṛe tunne}.

\begin{cousinweb}
\textbf{\pa{ਤੁੰਨਣਾ}} (\emph{tunnṇā}) is inherited from Sanskrit \emph{tunna}, \textquotedblleft{}struck\textquotedblright{}; CDIAL lists Punjabi \emph{tunnṇā}, \textquotedblleft{}to stow, stuff\textquotedblright{} (CDIAL 5863).
\end{cousinweb}

\begin{grammarlens}[title={What you've built}]
One more word for saying what you do.
\end{grammarlens}

\subsection*{Your turn}
\begin{itemize}
\raggedright
  \item \emph{Say it:} \emph{tunnṇā}
  \item \emph{Say it:} \emph{tunnṇā}, once more
  \item \emph{Say it:} \emph{maiṁ kappṛe baste vich tunndā hāṁ.} — I stuff my clothes into the bag.
  \item \emph{Recall:} say the Punjabi for memory, then the Punjabi for to come unstitched, then say \emph{tunnṇā} again
  \item \emph{Write it:} copy \textbf{\pa{ਤੁੰਨਣਾ}} once, checking the whole word
\end{itemize}

\begin{tcolorbox}[breakable,colback=teal!4,colframe=teal!35!black,title={Before you move on}]
What does \pa{ਤੁੰਨਣਾ} mean? (\textquotedblleft{}To stuff; to cram\textquotedblright{}.) Say it once more.
\end{tcolorbox}
